% sci_tier_access.tex — 4-tier learner-entry-point diagram for the SCI tutorial
% 4階層アクセス（学習者の入口）図。SCIチュートリアル用。
%
% Source of the four tiers: firmware/vehicle/docs/architecture.md §2
% 「学習者の入口（4階層アクセス）」の表を図として再構成したもの。
%
% Layout: four stacked boxes, top = highest abstraction (L0, beginners),
% bottom = lowest abstraction (L3, firmware implementers). A downward arrow
% chain to the right of the boxes shows that a learner may step down as far
% as their topic requires; each tier is otherwise usable on its own.
% レイアウト: 4段の箱を縦に積む。上ほど抽象度が高く（L0・初心者向け）、
% 下ほど抽象度が低い（L3・ファーム実装者向け）。箱の右側の下向き矢印は、
% 学習者が必要な分だけ下の層へ降りられることを示す（各層は単独でも使える）。
\documentclass[border=10pt]{standalone}
\usepackage{tikz}
\usetikzlibrary{arrows.meta, positioning, calc}
\usepackage{luatexja}
\usepackage[match]{luatexja-fontspec}
% Cross-platform font: Hiragino (macOS) → Yu Gothic (Windows) → Noto (Linux)
% クロスプラットフォーム: ヒラギノ (macOS) → 游ゴシック (Windows) → Noto (Linux)
\usepackage{ifthen}
\newboolean{jfontset}\setboolean{jfontset}{false}
\IfFontExistsTF{Hiragino Kaku Gothic ProN}{%
  \setmainjfont{Hiragino Kaku Gothic ProN}%
  \setsansjfont{Hiragino Kaku Gothic ProN}%
  \setboolean{jfontset}{true}}{}
\ifthenelse{\boolean{jfontset}}{}{%
  \IfFontExistsTF{Yu Gothic}{%
    \setmainjfont{Yu Gothic}%
    \setsansjfont{Yu Gothic}%
    \setboolean{jfontset}{true}}{}}
\ifthenelse{\boolean{jfontset}}{}{%
  \IfFontExistsTF{Noto Sans CJK JP}{%
    \setmainjfont{Noto Sans CJK JP}%
    \setsansjfont{Noto Sans CJK JP}}{}}

\begin{document}

\definecolor{l0color}{RGB}{0, 100, 180}
\definecolor{l1color}{RGB}{40, 140, 60}
\definecolor{l2color}{RGB}{120, 80, 150}
\definecolor{l3color}{RGB}{110, 110, 110}
\definecolor{gray1}{RGB}{120, 120, 120}

% Plain-LaTeX helper macros for the colored text lines inside each box
% (kept as ordinary commands, not tikz styles, since they are used inside
% node text rather than as tikzpicture options)
% 箱の中の色付きテキスト用マクロ（ノードのテキスト内で使うのでtikzの
% オプションではなく普通のLaTeXコマンドとして定義する）
\newcommand{\ttier}[2]{{\bfseries\large\color{#1}#2}}
\newcommand{\tns}[2]{{\ttfamily\normalsize\color{#1}#2}}
\newcommand{\tsub}[1]{{\small\color{black!55}#1}}

\begin{tikzpicture}[
    tier/.style={
      rectangle, rounded corners=4pt, draw=#1, line width=1.5pt,
      fill=#1!8, minimum width=92mm, minimum height=25mm,
      align=left, inner xsep=10mm, inner ysep=2mm
    },
    arr/.style={-{Stealth[length=3.2mm]}, line width=1.6pt, gray1},
  ]

  % Four tiers, top (L0, most abstract) to bottom (L3, closest to hardware)
  % 4層。上（L0・最も抽象的）から下（L3・ハードウェアに最も近い）まで
  \node[tier=l0color] (L0) at (0, 0) {%
    \begin{minipage}{86mm}\raggedright
      \ttier{l0color}{L0: Workshop API}\quad\tns{l0color}{ws::*}\\[3pt]
      \tsub{典型ユーザー: Workshop 受講者}\\[1pt]
      \tsub{setup() / loop\_400Hz(dt)、30+ 関数で完結}
    \end{minipage}%
  };
  \node[tier=l1color, below=12mm of L0] (L1) {%
    \begin{minipage}{86mm}\raggedright
      \ttier{l1color}{L1: Topic API}\quad\tns{l1color}{sf::api::*}\\[3pt]
      \tsub{典型ユーザー: 推定・制御学習者}\\[1pt]
      \tsub{Topic を subscribe/publish、IEstimator/IController 差替え}
    \end{minipage}%
  };
  \node[tier=l2color, below=12mm of L1] (L2) {%
    \begin{minipage}{86mm}\raggedright
      \ttier{l2color}{L2: HAL Direct}\quad\tns{l2color}{stampfly::*Wrapper}\\[3pt]
      \tsub{典型ユーザー: HW 学習者}\\[1pt]
      \tsub{SPI / I2C を直接呼ぶ（Topic を介さない経路）}
    \end{minipage}%
  };
  \node[tier=l3color, below=12mm of L2] (L3) {%
    \begin{minipage}{86mm}\raggedright
      \ttier{l3color}{L3: BSP Internal}\quad\tns{l3color}{sf::internal::board}\\[3pt]
      \tsub{典型ユーザー: ファーム実装者}\\[1pt]
      \tsub{bus handle、起動順序を直接扱う}
    \end{minipage}%
  };

  % Downward arrow chain, connecting each box directly at its center
  % (Lk.south -> L(k+1).north), so the arrowhead touches the box below.
  % Fix: the previous version drew the arrows 10mm to the right of the
  % boxes, floating in open space and touching nothing (T3 violation).
  % 下向き矢印の連鎖。各箱を中心で直結する（Lk.south -> L(k+1).north）。
  % 矢頭が下の箱に接する。修正: 以前は箱の右10mmにオフセットして描画
  % しており、何にも接続しない宙に浮いた矢印だった（T3違反）。
  \draw[arr] (L0.south) -- (L1.north);
  \draw[arr] (L1.south) -- (L2.north);
  \draw[arr] (L2.south) -- (L3.north);

  % Single side label for the whole arrow chain (vertically centered)
  % 矢印の連鎖全体に対する説明ラベル（縦方向の中央に1つだけ配置）
  \coordinate (midy) at ($(L1.south)!0.5!(L2.north)$);
  \node[rotate=270, font=\small, text=gray1, align=center,
        anchor=center] at ($(midy -| L0.east) + (26mm, 0)$)
    {必要な分だけ降りられる\\（各層は単独でも使える）};

\end{tikzpicture}
\end{document}
